% Wallets workstream: methodology. Code: src/albert/research/flb_wallets.py
Polymarket records every fill on chain with the wallet on each side, so we can follow identified
traders. This subsection defines the Polymarket return measures, the exact wallet accounting, the
wallet classes, and the tests of concentration, persistence, retail dynamics and cross-venue
pricing. Notation follows \cref{sec:method}: $\q$ is the price of the side bought, $\won$ is 1 if
that side pays out, $\fee$ is the fee per contract and $\ret=(\won-\q-\fee)/(\q+\fee)$ is the return
per dollar of holding to settlement.

\paragraph{Information cut-off.}
For each market we fix an information time $t^{\mathrm{info}}$ that does not depend on the
outcome. For a sports game it is the scheduled start (\texttt{gameStartTime}); otherwise it is the
scheduled end date. If a market closed before that schedule, the event happened early, and we set
$t^{\mathrm{info}}$ to 24 hours before the close. Unless stated otherwise, all calibration and
trade-level return measures use only fills and prices strictly before $t^{\mathrm{info}}$, so
in-play sports trading and post-outcome sweeps at 99.9\cents\ are excluded.

\paragraph{Calibration from prices at fixed leads.}
For a lead $L\in\{1\text{ h}, 24\text{ h}, 7\text{ d}\}$, the price of market $m$ is the last
trade at or before $t^{\mathrm{info}}_m-L$, expressed in units of outcome 0 (``YES''): a trade in
outcome 1 at price $p$ is a YES price of $1-p$. Each market contributes both sides, YES at $p$ and NO
at $1-p$. We report the mean of $\ret$ by side-price bucket, with weights equal to the design weight
of the market (``market-weighted'') or the design weight times market volume
(``volume-weighted''). The calibration slope $\beta$ is the coefficient of the logistic regression
\eqref{eq:logit} estimated on the outcome-0 side of each market (the other side duplicates it); with
$\beta>1$ prices are too close to $\tfrac12$, longshots too dear and favourites too cheap.

\paragraph{Trade-level returns.}
A buy fill of outcome $k$ at price $p$ has $\q=p$ and $\won=\ind{k \text{ wins}}$. Polymarket
charges takers $\fee=\theta_m\q(1-\q)$ per share on fee-enabled markets, where $\theta_m$ is the
market's published rate (mostly 0.03--0.05 by category, 0 for geopolitics); makers pay nothing. We use
the rate in each market's fee schedule, and $\fee=0$ for maker fills and for markets without fees,
which covers almost every market before March 2026 (\cref{app:wallets-fees}). Each table states whether returns are
gross ($\fee=0$) or net.
Fill-level means are either ``dollar-weighted'' (weight = notional $\times$ design weight) or
``fill-weighted'' (weight = design weight).

\paragraph{Exact wallet accounting.}
We express each fill in YES-equivalent units. Buying $s$ shares of outcome 0 at $p$ adds $+s$ to
exposure $e$ and $-sp$ to cash $c$; selling reverses both signs. Buying outcome 1 at $p$ pays
$s(1-Y)$ at settlement, so it adds $-s$ to $e$ and $+s(1-p)$ to $c$. The PnL of wallet $i$ in market $m$
is then
\begin{equation}\label{eq:wallets-pnl}
  \pi_{im}=\textstyle\sum_{j\in(i,m)} c_j+Y_m\sum_{j\in(i,m)} e_j-\sum_{j\in(i,m)}\fee_j s_j ,
\end{equation}
where $Y_m=1$ if outcome 0 wins. In these units a split (\$1 for one share of each outcome) or a
merge is PnL-neutral, so \eqref{eq:wallets-pnl} is exact from trades alone. Each fill appears once
for the taker and once for every maker it matched. In a fully observed fee-free market, therefore,
$\sum_i e_{im}=0$ and $\sum_i\pi_{im}=0$. We use this identity to check the data. Where fees apply,
Polymarket takes the taker's fee in shares on buys, and the identity holds up to those shares.

\paragraph{Wallet classes.}
Following \citet{rothschildsethi2016}, who group Intrade accounts by their trading behaviour, we classify
traders by observable conduct rather than by identity. Each fill is classified from the wallet's own
history in the sample strictly before the fill.
Let $n$ be its earlier fills, $\mu$ the share of those that were maker fills, $\tau$ the share of
the markets it had entered that it had already traded in both YES-equivalent directions, and
$\tilde v$ the median notional of its last 200 fills, including the current one. Classes are
assigned in this order:
\begin{enumerate}
  \item \textbf{Market maker (MM):} $n\geq 20$, $\mu\geq 0.7$ and $\tau\geq 0.5$.
  \item \textbf{Whale:} $\tilde v\geq\$1{,}000$.
  \item \textbf{Active retail:} $\tilde v\leq\$100$ and $n\geq 20$.
  \item \textbf{Retail:} $\tilde v\leq\$100$ and $n<20$.
  \item \textbf{Mid:} all other fills ($\$100<\tilde v<\$1{,}000$).
\end{enumerate}
Account age, which separates new from established retail, comes from the creation date
(\texttt{join\_date}) that Polymarket reports for each wallet. Age is measured at the fill, so it
is known at trade time. Leaderboard rank and lifetime counters are measured today, after the fact,
and are used only descriptively. The robustness grid moves one cut-off at a time (small size \$50
or \$250; whale size \$500 or \$5{,}000; active at 10 or 50 earlier fills; looser or stricter MM
criteria). \Cref{tab:wallets-robust} reports the results.

\paragraph{Who buys and who sells longshots.}
Besides the longshot share of each class's buying, we measure net longshot flow. A fill in a token
priced at most 10\cents\ counts as a longshot purchase or sale at its notional. A fill in a token
priced at least 90\cents\ counts as the opposite trade in the complementary longshot, at that
longshot's cost $s(1-p)$. Net flows sum to zero across all participants, so the classes that buy
longshots on net are matched by the classes that sell them. For positions opened by a buy before
$t^{\mathrm{info}}$, the exit value added
$(\text{proceeds}-\text{shares sold}\times\won)/\text{cost}$ compares the PnL actually realised
with that of holding every share to settlement.

\paragraph{Concentration and persistence.}
Wallet PnL sums \eqref{eq:wallets-pnl} over the sampled markets. We report the share of all gains (losses)
held by the top $x$\% of winning (losing) wallets, which is the Lorenz curve read from the top, and the
corresponding Gini coefficients. For persistence,
the sample is split by resolution date into H1 (January 2024 to May 2025) and H2 (June 2025 to
October 2026). Wallets with at least five markets and \$50 bought in each half are ranked on the
shrunk H1 return
\begin{equation}\label{eq:wallets-shrink}
  \tilde R_i^{\mathrm{H1}}=\frac{\sum_{m\in\mathrm{H1}}\pi_{im}}{\sum_{m\in\mathrm{H1}}B_{im}+K},
\end{equation}
where $B_{im}$ is the amount bought and $K$ is the median amount bought per wallet-half. They are
then evaluated on their raw H2 return, using the Spearman rank correlation with a wallet bootstrap and
the H2 return of the top minus the bottom H1 quintile. Two controls separate skill from luck and style:
\begin{itemize}
  \item \textbf{Edge versus entry price.} The wallet's mean buy return minus the mean return of all
    buys in the same price bucket and half. This removes persistence that comes only from buying
    favourites rather than longshots.
  \item \textbf{Luck placebo.} We redraw each market's outcome as
    $\tilde Y_m\sim\text{Bernoulli}(p_m^{\mathrm{close}})$, where $p_m^{\mathrm{close}}$ is the last
    price before $t^{\mathrm{info}}$, and recompute \eqref{eq:wallets-pnl} with the same trades.
    Under the placebo, prices are calibrated and no wallet can forecast outcomes beyond
    $p^{\mathrm{close}}$, but wallets keep the execution advantage of trading at better prices than
    the close.
\end{itemize}

\paragraph{Retail dynamics and herding.}
Account creation dates are looked up for a pseudo-random subset of sampled wallets: those whose address is
divisible by eight, in address order. Inflow is the weekly count of creations among them. Around five salient events
fixed in advance (the 2024 US presidential election, Super Bowls LIX and LX, the 2025 papal
conclave and the 2025 New York City mayoral election), we compare creations in the 14 days centred
on the event with a baseline four to eight weeks earlier. Herding of small wallets is measured in
one-hour windows with at least five small (at most \$100) taker buys by retail or active-retail
wallets, using the measure of
\citet{lakonishok1992}:
\begin{equation}\label{eq:wallets-lsv}
  H_{mt}=\Bigl|\tfrac{k_{mt}}{n_{mt}}-\bar p_m\Bigr|-\E\Bigl|\tfrac{K}{n_{mt}}-\bar p_m\Bigr|,\qquad K\sim\mathrm{Bin}(n_{mt},\bar p_m).
\end{equation}
Here $k_{mt}$ counts YES-side buys in the window and $\bar p_m$ is the market's average YES share of
small buys. $H>0$ means buys cluster on one side more than independent orders would. To test for
reversal, we regress the YES-price change over the next $h\in\{1,6,24\}$ hours on the window's
dollar imbalance $I_{mt}\in[-1,1]$ and on the previous hour's price change, using only windows that
end before $t^{\mathrm{info}}$.

\paragraph{Cross-venue comparison.}
For Kalshi game markets matched one-to-one to Polymarket outcomes (\cref{sec:data-wallets}), we
compare synchronous hourly prices at leads of 1--24 hours before the scheduled start. The Kalshi
price is the bid--ask midpoint of the hourly candle and the Polymarket price is the hourly price
history. The gap is $g=p^{\mathrm{PM}}-p^{\mathrm{K}}$, grouped by the two venues' average side
price. A gap exceeds fees if $|g|$ is larger than the sum of the two taker fees at the observed
prices, $0.07\,p^{\mathrm{K}}(1-p^{\mathrm{K}})$ on Kalshi and $\theta_m p^{\mathrm{PM}}(1-p^{\mathrm{PM}})$
on Polymarket. The comparison is descriptive.

\paragraph{Inference.}
Standard errors are clustered by event (Polymarket event id, or game for the cross-venue tests), as
in \cref{sec:method}. Population estimates use the design weights of the stratified sample.
